\documentclass[aspectratio=169]{beamer}
\usepackage{listings}
\usepackage{xcolor}

\definecolor{codegray}{rgb}{0.5,0.5,0.5}
\definecolor{codepurple}{rgb}{0.58,0,0.82}
\definecolor{backcolour}{rgb}{0.95,0.95,0.92}

\lstdefinestyle{customc}{
  backgroundcolor=\color{backcolour},
  commentstyle=\color{codegray}\itshape,
  keywordstyle=\color{blue}\bfseries,
  numberstyle=\tiny\color{gray},
  stringstyle=\color{codepurple},
  basicstyle=\ttfamily\footnotesize,
  breakatwhitespace=false,
  breaklines=true,
  captionpos=b,
  keepspaces=true,
  numbers=left,
  numbersep=5pt,
  showspaces=false,
  showstringspaces=false,
  showtabs=false,
  tabsize=2,
  frame=single,
  language=C++
}

\lstset{style=customc}

\usepackage[spanish]{babel}
\usepackage[utf8]{inputenc}
\usepackage[T1]{fontenc}
\usepackage{lmodern}
\usepackage{amsmath, amsfonts}
\usepackage{hyperref}

\usetheme{Madrid}

\title{Costo asintótico}
\author{Sebastian Mestre}
\date{}

\begin{document}

\begin{frame}
  \titlepage
\end{frame}

\section{Motivación}

\begin{frame}{¿Por qué nos importa el costo?}
  \begin{itemize}
    \item En programación competitiva nos exigen que un programa termine en \textbf{menos de cierto tiempo} (típicamente 1 segundo).
    \item Para saber si nuestro programa va a entrar en tiempo, necesitamos estimar \textbf{cuánto tarda} antes de implementarlo completo.
    \item Implementar una idea y recién después descubrir que es demasiado lenta significa \textbf{perder mucho tiempo} en concurso.
    \item Queremos poder razonar sobre el \textbf{tiempo de ejecución} solo pensando en el algoritmo.
  \end{itemize}
\end{frame}

\begin{frame}{Tiempo exacto vs aproximación}
  \begin{itemize}
    \item La cantidad exacta de tiempo que tarda un programa es imposible de conocer de antemano:
      \begin{itemize}
        \item Depende de la máquina, compilador, caché, sistema operativo, qué otro software está corriendo, etc.
      \end{itemize}
    \item Datos típicos:
      \begin{itemize}
        \item Límite de tiempo \(\approx 1\) segundo.
        \item Una CPU moderna ejecuta del orden de \(3 \cdot 10^9\) operaciones por segundo.
      \end{itemize}
    \item \textbf{Idea:} alcanza con garantizar que nuestro programa hace \textbf{menos de \(3 \cdot 10^9\)} operaciones.
  \end{itemize}
\end{frame}

\section{Modelo y notación O}

\begin{frame}{Dependencia en el tamaño de la entrada}
  \begin{itemize}
    \item El tiempo de ejecución depende del \textbf{tamaño de la entrada}, denotado por \(N\)
    \item \textbf{Idea:} buscamos una \textbf{cota superior} del número de operaciones en función de \(N\).
    \item Como \(3 \cdot 10^9\) es un número grande, podemos permitirnos \textbf{aproximaciones asintóticas} (qué pasa cuando \(N \to \infty\)).
    \item Ejemplo: cuando \(N\) es grande
      \[
        N^2 + N \sim N^2.
      \]
%  \end{itemize}
%\end{frame}
%
%\begin{frame}{Solo importa el orden de magnitud}
%  \begin{itemize}
    \item En práctica competitiva el límite de tiempo suele ser bastante holgado, incluso \textbf{20x más de lo necesario}.
      \begin{itemize}
        \item Nos importa diferenciar \(N\) de \(N^2\).
        \item Pero no distinguir entre \(0{,}5N\) y \(10N\).
      \end{itemize}
    \item Para esto usamos la \textbf{notación O grande}.
  \end{itemize}
\end{frame}

\begin{frame}{Definición formal de notación O}
  \begin{block}{\textbf{Notación O}}
    Matemáticamente, si \(f : \mathbb{N} \rightarrow \mathbb{R}^+\) representa el tiempo de ejecución en función del tamaño de la entrada \(N\), entonces:
    \begin{itemize}
      \item \(O(f(N))\) es el conjunto de todas las funciones \(g : \mathbb{N} \rightarrow \mathbb{R}^+\) tales que
      \item existen una constante \(C > 0\) y un natural \(N_0\) tal que para todo \(N \ge N_0\),
        \[
          g(N) \le C \cdot f(N).
        \]
    \end{itemize}
  \end{block}
  
  Cuando decimos que un programa tiene \textbf{costo} \(O(f)\), significa que la función que representa el \textbf{peor caso} de tiempo de ejecución para entradas de tamaño \(N\) pertenece a \(O(f)\).
\end{frame}

\begin{frame}{Definición informal de notación O}
\begin{figure}
    \centering
    \includegraphics[width=0.5\linewidth]{image.png}
    
    Una funcion es $O(f)$ si, para valores grandes de N, sus picos mas altos no superan a f multiplicada por alguna constante.
\end{figure}
\end{frame}

\section{Reglas prácticas}

\begin{frame}[fragile]{Operaciones básicas: \(O(1)\)}
  \begin{itemize}
    \item En el modelo usual, muchas operaciones simples tienen \textbf{costo constante} \(O(1)\).
  \end{itemize}
  \begin{block}{Ejemplo}
  \begin{lstlisting}[style=customc]
int x = 0;
x += 1; // O(1)
x *= 2; // O(1)
x %= 3; // O(1)
  \end{lstlisting}
  \end{block}
\end{frame}

\begin{frame}[fragile]{Bucles lineales: \(O(N)\)}
  \begin{itemize}
    \item Un bucle típico de \(N\) pasos tiene costo \(O(N)\).
  \end{itemize}
  \begin{block}{Ejemplo}
  \begin{lstlisting}[style=customc]
int sum = 0;
for (int i = 0; i < N; ++i) { // O(N) iteraciones
  sum += i; // O(1)
}
// total: O(N)
  \end{lstlisting}
  \end{block}
\end{frame}

\begin{frame}[fragile]{Bucles anidados: multiplicar costos}
  \begin{itemize}
    \item Si tenemos bucles anidados, los costos se \textbf{multiplican}.
  \end{itemize}
  \begin{block}{Ejemplo}
  \begin{lstlisting}[style=customc]
for (int i = 0; i < N; ++i) {   // O(N) iteraciones
  for (int j = 0; j < M; ++j) { // O(M)
    ...
  }
}
// total: O(N * M)
  \end{lstlisting}
  \end{block}
\end{frame}

\begin{frame}[fragile]{Sumar costos: solo el término dominante}
  \begin{itemize}
    \item Código en líneas consecutivas se \textbf{suma}.
    \item Pero en notación O solo nos importa el \textbf{término superior}.
  \end{itemize}
  \begin{block}{Ejemplo}
  \begin{lstlisting}[style=customc]
for (int i = 0; i < N*N; ++i) { ... }     // O(N^2)
for (int i = 0; i < 3*N+15; ++i) { ... }  // O(N)
for (int i = 0; i < N*N-5; ++i) { ... }   // O(N^2)
// total: O(N^2 + N + N^2) = O(N^2)
  \end{lstlisting}
  \end{block}
\end{frame}

\section{Patrones típicos}

\begin{frame}[fragile]{Algoritmos de dos punteros}
  \begin{itemize}
    \item Hay patrones que \textbf{rompen} la intuición simple de “bucle dentro de bucle = \(O(N^2)\)”.
    \item Un ejemplo clásico son los \textbf{algoritmos de dos punteros}.
  \end{itemize}
  \begin{block}{Ejemplo}
  \begin{lstlisting}[style=customc]
int j = 0;
for (int i = N; i > 0; --i) {             // O(N) iteraciones
  while (j < N && a[j] + a[i-1] < x) {    // O(N)
    j++;
  }
}
// total: O(N)
// aunque a primera vista parezca O(N^2)
  \end{lstlisting}
  \end{block}
\end{frame}

\begin{frame}[fragile]{Patrones logarítmicos}
  \begin{itemize}
    \item Cuando un número se va \textbf{multiplicando o dividiendo por 2} hasta alcanzar una cota, suele aparecer un \(\log N\).
  \end{itemize}
  \begin{block}{Ejemplo}
  \begin{lstlisting}[style=customc]
for (int x = 1; x < N; x *= 2) { ... }
for (int x = N; x > 1; x /= 2) { ... }
// total: O(log N)
  \end{lstlisting}
  \end{block}
\end{frame}

\begin{frame}[fragile]{Suma armónica: \(O(N \log N)\)}
  \begin{itemize}
    \item Aparece a veces (e.g. Criba de Eratostenes, algunas DPs) un patrón donde el número de iteraciones del bucle interno es \(N / i\).
    \item La suma resultante se comporta como \(N \log N\).
  \end{itemize}
  \begin{block}{Ejemplo}
  \begin{lstlisting}[style=customc]
for (int i = 1; i < N; ++i) {         // O(N) iteraciones
  for (int j = i; j < N; j += i) {    // O(N / i) iteraciones
    ... // O(1)
  }
}
// total: O(N/1 + N/2 + ... + N/N) = O(N log N)
  \end{lstlisting}
  \end{block}
\end{frame}

\section{Regla del pulgar para concursos}

\begin{frame}{¿Qué costo es aceptable para cada \(N\)?}
  \begin{itemize}
    \item Asumimos límite de tiempo \(\approx 1\) segundo y una implementación razonable en C++.
  \end{itemize}
  \vspace{0.3cm}
  \begin{center}
  \begin{tabular}{|c|c|}
    \hline
    \textbf{N} & \textbf{Costo} \\
    \hline
    \(\le 100\)   & \(O(N^4)\) \\
    \(\le 500\)   & \(O(N^3)\) \\
    \(\le 10^4\)  & \(O(N^2)\) \\
    \(\le 10^5\)  & \(O(N \sqrt{N})\) \\
    \(\le 10^6\)  & \(O(N \log N)\) \\
    \(\le 10^7\)  & \(O(N)\) \\
    \(\le 10^{18}\) & \(O(\log N)\) \\
    \hline
  \end{tabular}
  \end{center}
\end{frame}

\section*{Cierre}

\begin{frame}{Ideas clave}
  \begin{itemize}
    \item No podemos conocer el \textbf{tiempo exacto}, pero sí su \textbf{orden de magnitud}, y con eso alcanza.
    \item La notación \(O(\cdot)\) nos permite razonar sobre el \textbf{peor caso} de tiempo en función de \(N\).
    \item Aprender a reconocer patrones típicos (\(O(N)\), \(O(N \log N)\), \(O(N^2)\), etc.) es crucial para programar bajo límite de tiempo.
  \end{itemize}
\end{frame}

\begin{frame}{Preguntas}
  \centering
  \Huge
  \textbf{Q\&A}
\end{frame}

\end{document}
